\chapter*{第二册能力清单}
\addcontentsline{toc}{chapter}{第二册能力清单}

第二册的出口不是“知道很多系统名词”，而是能把性能、并发、运行时、I/O 和分布式边界写成可验证的工程报告。

\section*{能实现}

\begin{itemize}
  \item 实现顺序 reference、并行 reduce/scan/partition、线程切分和局部 partial 合并。
  \item 实现可关闭 bounded queue，覆盖阻塞 \code{push}、阻塞 \code{pop}、\code{close}、drain、等待报告和容量为一压力测试。
  \item 实现一个最小 thread pool 或 task runtime，能解释 ready、running、blocked、finished、failed、canceled 和 \code{wait\_idle}。
  \item 实现 I/O 管线的最小状态机：buffer ownership、short write、completion、manifest、cancel 和 backpressure。
\end{itemize}

\section*{能诊断}

\begin{itemize}
  \item 用 FLOPs、bytes、arithmetic intensity、工作集、页数和线程扩展曲线判断 kernel 可能受什么限制。
  \item 用 PMU/Top-Down 或替代证据区分 frontend bound、bad speculation、backend bound、memory bound、cache/TLB miss 和 branch miss。
  \item 诊断 false sharing、lock convoy、futex 等待、same-pool future 饥饿、barrier phase skip 和 broken promise。
  \item 区分 eventfd/epoll readiness、io\_uring completion、page cache、direct I/O、fsync 和可靠提交。
\end{itemize}

\section*{能解释}

\begin{itemize}
  \item 解释 cache line、TLB、page size、huge page、NUMA first-touch、coherence、HITM 和 remote access 的证据边界。
  \item 解释 AVX2/AVX-512/AMX 类后端优化为什么要看 lane、tail、horizontal reduction、gather/scatter、port pressure 和 downclock。
  \item 解释 C++ memory model 中 data race、happens-before、release/acquire、relaxed、seq\_cst 和 litmus test。
  \item 解释 lock-free 的 linearization point、ABA、hazard pointer、epoch reclamation、RCU 和 progress guarantee。
\end{itemize}

\section*{不能夸大的部分}

第二册主要训练现代单机计算系统和分布式边界的基本方法。它不替代第三册的 tokenizer、真实模型加载、KV cache、quantization、AI compiler、serving 和 GPU kernel 训练。
